\clearpage
\Article{348}
\textsc{Aufgabe.} \emph{Gegeben sei die Gleichung}
\[
x^m+Ax^{m-1}+Bx^{m-2}+Cx^{m-3}+\text{usw.}+M=0
\]
\emph{mit den Wurzeln $x=a$, $x=b$, $x=c$ usw. Man finde eine Gleichung mit den Wurzeln $x=a^\tau$, $x=b^\tau$, $x=c^\tau$ usw.}

\SourceLeaf{233}
\emph{Erste Lösung.} Man bestimme mit dem bekannten Satz die Summen der Wurzeln der gegebenen Gleichung, ihrer Quadrate, Kuben usw. bis zur Potenz mit dem Exponenten $m\tau$. Damit erhält man auch die Summen der Wurzeln der gesuchten Gleichung und ihrer Quadrate usw., nämlich $\Sigma a^\tau$, $\Sigma a^{2\tau}$ usw. Hieraus lassen sich mit demselben Satz die Koeffizienten bestimmen.

Für die praktische Rechnung ist diese Lösung zwar einfacher; doch eignet sich die folgende besser für unseren Zweck und auch dazu, zu zeigen, dass die Koeffizienten der gesuchten Gleichung ganze Zahlen sind, wenn die Koeffizienten der gegebenen Gleichung ganze Zahlen sind.

\Needspace{23\baselineskip}
\emph{Zweite Lösung.} Es sei $\theta$ eine primitive Wurzel der Gleichung $x^\tau=1$. Man bilde das Produkt aus
\[
\begin{aligned}
&x^m+Ax^{m-1}+Bx^{m-2}+\text{usw.}\\
&x^m+A\theta x^{m-1}+B\theta\theta x^{m-2}+\text{usw.}\\
&x^m+A\theta\theta x^{m-1}+B\theta^4x^{m-2}+\text{usw.}\\
&\text{usw.}\\
&x^m+A\theta^{\tau-1}x^{m-1}+B\theta^{2\tau-2}x^{m-2}+\text{usw.}
\end{aligned}
\]
Die Wurzeln dieses Produkts sind daher
\[
\begin{array}{llll}
a,&\theta a,&\theta\theta a&\text{usw.}\\
b,&\theta b,&\theta\theta b&\text{usw.}\\
c,&\theta c,&\theta\theta c&\text{usw.}
\end{array}
\]
d. h. das Produkt ist gleich
\[
(x^\tau-a^\tau)(x^\tau-b^\tau)(x^\tau-c^\tau)\ldots
\]
und hat somit die Form
\[
x^{\tau m}+A'x^{\tau(m-1)}+B'x^{\tau(m-2)}+\text{usw.}
\]
\Needspace{14\baselineskip}
Schreibt man nun für $x^\tau$ die Größe $x$, so ist
\[
x^m+A'x^{m-1}+B'x^{m-2}+\text{usw.}=(x-a^\tau)(x-b^\tau)(x-c^\tau)\ldots
\]
und damit
\[
x^m+A'x^{m-1}+B'x^{m-2}+\text{usw.}=0
\]
die gesuchte Gleichung. Dass hier $A',B'$ usw. nicht nur rationale, sondern auch ganze Zahlen sind, folgt leicht aus der Theorie der Gleichung $x^\tau=1$.

Da wir diese Operation im Folgenden oft verwenden werden, bezeichnen wir mit $(P,\rho^\tau)$
\SourceLeaf{234}
die Funktion, die, gleich null gesetzt, eine Gleichung ergibt, deren Wurzeln die Potenzen mit dem Exponenten $\tau$ der Wurzeln der Gleichung $P=0$ sind.

Ist $P\equiv Q$ nach irgendeinem Modul, so ist auch $(P,\rho^\tau)\equiv(Q,\rho^\tau)$ nach demselben Modul.

\Needspace{10\baselineskip}
\Article{349}
\textsc{Satz.} \emph{Der Koeffizient des Terms $x^n$ in $(P,\rho^\tau)$ ist nach dem Modul $\tau$ dem Koeffizienten des Terms $x^{\tau n}$ in $P^\tau$ kongruent, sofern $\tau$ eine Primzahl ist} (was für diesen Fall eine dritte Lösung der vorigen Aufgabe ergibt).

\emph{Beweis.} Aus dem sechsten Kapitel folgt, dass jeder Koeffizient des Produkts
\[
(x^m+Ax^{m-1}+\text{usw.})(x^m+A\theta x^{m-1}+\text{usw.})\ldots
\]
nach dem Ersetzen von $\theta^\tau$ durch die Einheit die folgende Form hat:
\[
E+(1+\theta+\theta\theta+\text{usw.}+\theta^{\tau-1})F
\]
Betrachtet man nun $\theta$ als primitive Wurzel der Gleichung $x^\tau=1$, so geht das ganze Produkt in $E$ über. Setzt man hingegen $\theta=1$, so geht das ganze Produkt in $P^\tau=E+\tau F$ über. Daher ist der Koeffizient des Terms $x^{n\tau}$ in $P^\tau$ nach dem Modul $\tau$ dem Koeffizienten des Terms $x^{n\tau}$ in $E$ kongruent, d. h. dem Koeffizienten des Terms $x^n$ in $(P,\rho^\tau)$.

\Needspace{10\baselineskip}
\Article{350}
\textsc{Satz.} \emph{Ist $\tau$ eine Primzahl, so ist}
\[
(P,\rho^\tau)\equiv P\;(\mathit{mod.}\,\tau)
\]

\begingroup\postdisplaypenalty=10000\relax \emph{Beweis.} Es sei der Koeffizient des Terms $x^n$ in $(P,\rho^\tau)$ $=N'$, der Koeffizient desselben Terms in $P$ dagegen $=N$. Setzt man dann
\[
P=x^m+Ax^{m-1}+\text{usw.}+Nx^n+\text{usw.}
\]
so gilt
\[
P^\tau\equiv x^{m\tau}+A^\tau x^{(m-1)\tau}+\text{usw.}+N^\tau x^{n\tau}+\text{usw.}\;(\mathrm{mod.}\,\tau)
\]
und daher nach dem vorigen Artikel $N'\equiv N^\tau\;(\mathrm{mod.}\,\tau)$. Wegen $N^\tau\equiv N$ ist somit $N'\equiv N$. Q.\,E.\,D.\par\endgroup

Hieraus ist auch ersichtlich, dass $(P,\rho^\alpha)\equiv(P,\rho^{\alpha\tau})$ und $(P,\rho^\tau)\equiv(P,\rho^{\tau\tau})$, also allgemein
\[
(P,\rho^\alpha)\equiv(P,\rho^{\alpha\tau^k})\;(\mathrm{mod.}\,\tau)
\]

\SourceLeaf{235}
\Needspace{10\baselineskip}
\Article{351}
\textsc{Satz.} \emph{Es gibt einen Wert der Zahl $\nu$, der kleiner als $p^m$ ist, sodass die Funktion $x^\nu-1$ durch die gegebene Funktion $P$ vom Grad $m$, deren niedrigster Term die Unbestimmte $x$ nicht enthält, nach dem Modul $p$ teilbar ist.}

\emph{Beweis.} Man teile durch $P$ die Folge der Funktionen $1,x,xx\ldots$ bis $x^{p^m-1}$, sobald ihr Grad $m$ übersteigt. Da keine durch $P$ ohne Rest teilbar ist, lassen sich alle Reste auf die Form
\[
Ax^{m-1}+Bx^{m-2}+\dotstwo+N
\]
bringen, wobei alle Koeffizienten positiv und $<p$ sind. Da aber niemals alle $=0$ sein können, gibt es offenbar nur $p^m-1$ Funktionen, deren einer die einzelnen Reste gleich sein müssen. Da man bis zur Potenz von $x$ mit dem Exponenten $p^m-1$ insgesamt $p^m$ Reste erhält, müssen notwendig mindestens zwei übereinstimmen. Bei der Division von $x^a$ und $x^{a+\nu}$ durch $P$ möge also derselbe Rest entstehen, wobei $a+\nu<p^m$. Dann ist $x^{a+\nu}-x^a$ durch $P$ teilbar. Da nach Voraussetzung $x$ und damit auch $x^a$ eine zu $P$ teilerfremde Funktion ist, lässt sich daher auch $x^\nu-1$ durch $P$ teilen. Q.\,E.\,D.

\emph{Folgerung.} Ist $x^\nu-1$ durch $P$ teilbar, so ist auch $x^{k\nu}-1$ durch $P$ teilbar, wobei $k$ eine beliebige ganze Zahl bezeichnet.

\Needspace{10\baselineskip}
\Article{352}
\textsc{Satz.} \emph{Bei den Bezeichnungen des vorigen Artikels sei $P$ eine Primfunktion und $x^\nu$ die niedrigste Potenz, die nach Abzug der Einheit durch $P$ teilbar ist. Dann ist $\nu$ entweder $=p^m-1$ oder ein Teiler dieser Zahl, mit der einzigen Ausnahme $P=x$.}

\Needspace{26\baselineskip}
\emph{Beweis.} Da $P$ eine Primfunktion vom Grad $m$ ist, gibt es $p^m-1$ verschiedene Funktionen von kleinerem Grad als $m$, wenn man die Funktion $0$ aus ihrer Gesamtzahl ausschließt. Sie sind alle zu $P$ teilerfremd. Da nun $x^\nu$ der Annahme nach die niedrigste Potenz ist, die bei der Division durch $P$ die Einheit als Rest lässt, entstehen offenbar bei der Division aller niedrigeren Potenzen von $1,x,\dotstwo$ bis $x^{\nu-1}$ durch $P$ genau $\nu$ verschiedene Reste, die allgemein mit $A$ bezeichnet seien. Erschöpfen sie alle möglichen Reste, so ist der Satz bewiesen. Sind aber noch einige nicht darunter, so sei $B$ irgendeiner davon. Offenbar ergibt die Funktion $Bx^\nu$ bei der Division durch $P$ den Rest $B$, und allgemein ist $Bx^{\nu+k}\equiv Bx^k(\mathrm{mod.}\,P)$. Alle Funktionen von $B$ bis $Bx^{\nu-1}$ liefern aber
\SourceLeaf{236}
Reste, die untereinander und von den Resten $A$ verschieden sind. Wäre nämlich $Bx^\lambda\equiv Bx^{\lambda+\delta}(\mathrm{mod.}\,P)$, so wäre auch $1\equiv x^\delta(\mathrm{mod.}\,P)$ und $\delta<\nu$, entgegen der Voraussetzung. Wäre dagegen $Bx^\lambda\equiv x^\mu(\mathrm{mod.}\,P)$, so wäre $B\equiv x^{\mu+\nu-\lambda}(\mathrm{mod.}\,P)$ und damit $B$ einer der Reste $A$, ebenfalls entgegen der Voraussetzung. Man hat somit offenbar noch $\nu$ neue Reste. Auf dieselbe Weise kann man weiter fortfahren (ganz wie oben, \textsection.\ .). Dabei zeigt sich, dass die Anzahl $p^m-1$ aller möglichen Reste entweder $=\nu$ oder $=2\nu$ oder $=3\nu$ oder allgemein ein Vielfaches der Zahl $\nu$ ist. Q. E. D.

\Needspace{10\baselineskip}
\Article{353}
Aus dem vorigen Satz und der Folgerung zu \textsection.\ 351 ergibt sich, dass jede Primfunktion vom Grad $n$ die Funktion $x^{p^n-1}-1$ nach dem Modul $p$ teilt. Daher teilen alle Funktionen ersten Grades mit der einzigen Ausnahme derjenigen, die $\equiv x$ ist, die Funktion $x^{p-1}-1$. Dies ist der \textsc{Fermatsche Satz}. Alle Primfunktionen zweiten Grades, d. h. von der Form $xx+Ax+B$, teilen die Funktion $x^{pp-1}-1$ usw. Die Zahl $n$ habe nun sämtliche Teiler $n,\delta,\delta',\delta''$ usw.\ .\ $1$. Offenbar ist $p^n-1$ auch durch $p^\delta-1$, $p^{\delta'}-1$, $p^{\delta''}-1$ usw. $p-1$ teilbar. Deshalb ist die Funktion $x^{p^n-1}-1$ durch \emph{alle Primfunktionen der Grade $n,\delta,\delta',\delta''$ usw. bis zu den Primfunktionen ersten Grades (unter Ausschluss der Funktion $x$)} teilbar und daher auch, da alle diese Funktionen an sich prim und somit untereinander teilerfremd sind, durch \emph{das Produkt aus allen}. Eben dieses Produkt hat aber den Grad $p^n-1$ (\textsection.\ 347), da die eine Funktion $x$ fehlt. Somit ist klar, dass \emph{dieses Produkt selbst $x^{p^n-1}-1\;(\mathrm{mod.}\,p)$ kongruent sein muss.}

\Needspace{20\baselineskip}
\Article{354}
\textsc{Satz.} \emph{Ist die Funktion $x^\nu-1$ durch die Funktion $P$ teilbar, so gilt}
\[
(P,\rho^{k\nu+t})\equiv(P,\rho^t)
\]
\emph{wobei $k,t$ beliebige ganze Zahlen bezeichnen.}

\emph{Beweis.} Es sei
\[
P=x^m+Ax^{m-1}+Bx^{m-2}+\text{usw.}
\]
Bekanntlich ergibt die Entwicklung von
\[
\frac{mx^{m-1}+(m-1)Ax^{m-2}+\text{usw.}}{x^m+Ax^{m-1}+\text{usw.}}
\]
in eine unendliche Reihe der Form
\[
m\frac1x+\alpha\frac1{xx}+\beta\frac1{x^3}+\gamma\frac1{x^4}+\text{usw.}
\]
\SourceLeaf{237}
für $\alpha$ die Summe der Wurzeln der Gleichung $P=0$, für $\beta$ die Summe ihrer Quadrate usw. Hieraus folgt ohne Mühe, dass die Summen der Potenzen mit den Exponenten $\nu+1$, $\nu+2$ usw. der Summe der Wurzeln, ihrer Quadrate usw. kongruent sind. Sofern der Modul nicht gleich dem Grad der Funktion $P$ oder kleiner ist, folgt hieraus
\[
(P,\rho^{\nu+1})\equiv P,\quad(P,\rho^{\nu+2})\equiv(P,\rho^2),\quad(P,\rho^{\nu+3})\equiv(P,\rho^3)\text{ usw.}
\]
Jenen Fall werden wir weiter unten betrachten.

\Needspace{20\baselineskip}
\Article{355}
\textsc{Satz.} \emph{Sind in der Folge}
\[
(P,\rho^0),\quad(P,\rho)\quad(P,\rho^2),\quad(P,\rho^3)\text{ usw.}
\]
\emph{die auf das Glied an der Stelle $\nu$ folgenden Glieder den ersten der Reihe nach kongruent, so ist $x^\nu-1$ durch $P$ teilbar, sofern $P$ keinen Faktor mehrfach enthält.}

\emph{Beweis.} Setzt man $\frac{\mathrm{d}P}{\mathrm{d}x}=Q$, so ist $Q$ eine zu $P$ teilerfremde Funktion. Es sei
\[
\frac QP\equiv\frac Ax+\frac B{xx}+\frac C{x^3}+\text{usw.}
\]
Dann folgt nach Voraussetzung auf das Glied $\frac N{x^\nu}$
\[
\frac A{x^{\nu+1}}+\frac B{x^{\nu+2}}+\frac C{x^{\nu+3}}+\text{usw.}
\]
Daher ist
\[
\frac QP\equiv\frac{Ax^{\nu-1}+Bx^{\nu-2}+\text{usw.}}{x^\nu-1}
\]
woraus sich ergibt, dass die Funktion $x^\nu-1$ durch $P$ teilbar ist. Q. E. D.
